\chapter{Invariant Computation}
\label{chap:invariant-computation}

% Closed question: What remains invariant under admissible relabeling or
% recompilation?

\begin{chapteressay}

The previous chapter ended with the commutator: the algebraic record
of the failure of operations to commute. The chapter showed that
when two operations applied in different orders produce different
results, the difference is a structured quantity. This chapter asks
the dual question: what does not change? When two operations are
applied in different orders, what survives the relabeling?

The answer is the invariants. Some quantities depend on the
operations and the order; some do not. The quantities that do not
depend on the order, or the labeling, or the specific representation
are the invariants. The invariants are the algorithmic content that
survives admissible transformation.

The chapter's first concrete example is bootstrapping a compiler.
The Lean 4 compiler is written in Lean 4. Compiling the compiler
with itself should produce a semantically equivalent compiler.
The Lean community has explicitly demonstrated this:
\texttt{lean-stage-1} (the compiler produced by Stage 0 compiling
the Lean 4 source) and \texttt{lean-stage-2} (the compiler
produced by \texttt{lean-stage-1} compiling the same source) are
equivalent under the chosen build equivalence --- the binaries
agree on what they compute, modulo timestamps, layout, and other
build metadata. The bootstrap has converged to a fixed point up
to that equivalence.

The fixed point is the invariant. The compiler's binary is
what survives the recompilation. The recompilation is the
admissible relabeling. The invariance is the proof that the
algorithm has reached self-consistency.

The chapter develops invariance through five algorithmic forms.
Bootstrapping a compiler (Medium section, the load-bearing
example). Gradient descent invariance (neural network training
is symmetric under hidden neuron relabeling). Ising model phase
transition (the order parameter is the invariant of the broken
$\mathbb{Z}_2$ symmetry). Feynman diagram as computation
(amplitudes are invariant under internal momentum relabeling).
\texttt{ElaborationProcess} in Lean (the typeclass that captures
the stages of compilation, with the \texttt{HALTED} certificate
marking the final invariant).

The second concrete example is the theater production. A play has
been running for fifteen years. The cast has rotated entirely,
the set has been rebuilt twice, the lighting rig has been
replaced, three directors have come and gone. And yet the play is
``the same play.'' What survives is the relational structure: the
script, the blocking, the timing, the emotional arc. What changes
is the implementation: the specific actors, the specific
scenery, the specific lighting cues.

The play is the type. The production is the binary. Different
binaries implement the same type. The implementations differ in
details; the type is invariant. The director's job is to maintain
the invariance: when bringing on new cast members, the director
shapes the performance to remain within the play's equivalence
class.

The chapter's closed question is:

\begin{center}
\emph{What remains invariant under admissible relabeling or
recompilation?}
\end{center}

The answer the chapter develops is: the type-level content. The
internal details of the algorithm can vary across runs; the
type-level result is the same. The compiler's binary is the
invariant; the elaboration trace is the variable form. The
neural network's loss is the invariant; the specific weights are
the variable form. The play's structure is the invariant; the
specific cast is the variable form.

The chapter's algorithmic claim is that invariance is what
algorithmic discipline produces. The discipline of typing,
elaboration, type-checking, and bootstrapping is the
algorithmic procedure that produces the invariant content. The
content is the result of the type-check, the fixed point of the
bootstrap, the order parameter of the symmetry-broken phase, the
amplitude of the Feynman diagram, the play's emotional arc.

The chapter is short, but its content is structurally important.
Invariance is what makes computation trustworthy. The compiler
can be trusted because its output is invariant under
recompilation. The trained network can be trusted because its
loss is invariant under hidden neuron relabeling. The play can
be trusted to be the same play because its essential structure
is invariant under cast rotation. Each invariance is the
algorithm's commitment that, despite the variable details, the
essential content is preserved.

The book's final chapter takes up the append-only structure of
the ledger and the type-check that closes the book. The
invariants of this chapter are the foundation for the closing
discipline of the next.

\end{chapteressay}

\section{Bootstrapping a Compiler}
\label{se:bootstrapping-a-compiler}

The Lean 4 compiler is itself written in Lean 4. Before Lean 4 existed,
this was a chicken-and-egg problem: to compile Lean 4 source, you
need a Lean 4 compiler; to write a Lean 4 compiler, you need to be
able to write code in some language. The resolution is bootstrapping.

The bootstrap procedure works as follows. Stage 0 is a Lean 4
compiler written in some other language, typically C++ (Lean's
ancestor Lean 3 was a C++/Lean hybrid). Stage 0 cannot be written
in Lean 4 itself because no Lean 4 compiler exists yet. Stage 0
is hand-coded, debugged, and tested until it can compile some
significant subset of Lean 4 source.

Stage 1 is the Lean 4 source of a Lean 4 compiler. The source is
written by hand, in Lean 4. The source is not yet executable: it is
source code. Compile it using Stage 0. The output is a binary that
implements the Lean 4 compiler in Lean 4. Call this binary
\texttt{lean-stage-1}.

Now run \texttt{lean-stage-1} on the same Lean 4 source. The output
is a binary that implements the Lean 4 compiler in Lean 4. Call
this binary \texttt{lean-stage-2}.

The question is: are \texttt{lean-stage-1} and \texttt{lean-stage-2}
the same binary? If yes, the compiler has reached a fixed point.
If no, something is wrong: \texttt{lean-stage-2} is the Lean 4
compiler compiled by the Lean 4 compiler, and any discrepancy
between it and \texttt{lean-stage-1} indicates that Stage 0 was
producing different code than Lean would have produced.

For the Lean 4 bootstrap, the fixed point is achieved: the binary
output by \texttt{lean-stage-1} and the binary output by
\texttt{lean-stage-2} are semantically equivalent --- the same
inputs produce the same outputs --- modulo timestamps, layout,
and other build metadata. The compiler has converged up to the
chosen build equivalence.

This is the chapter's structure in its most precise form. The
compiler is an algorithm. Recompiling the compiler with itself
should produce the same compiler. The same compiler is the
invariant: it is what survives recompilation. The recompilation
process is the admissible relabeling. The invariance is the proof
that the algorithm has reached a fixed point.

Why does this matter? Because it certifies that the compiler is
internally consistent. If \texttt{lean-stage-1} and
\texttt{lean-stage-2} differed, the discrepancy would indicate that
the compiler's behavior depends on the underlying compiler (Stage
0), not solely on the Lean 4 specification. The fixed-point
property says: the compiler's behavior depends only on the Lean 4
specification. The underlying compiler used to bootstrap is
incidental.

The fixed-point property is also operationally important. It means
that, in principle, the Lean 4 community could discard Stage 0
entirely once \texttt{lean-stage-1} exists. \texttt{lean-stage-1}
can recompile itself; subsequent versions can be produced by
running the current version on the updated source. The Stage 0
compiler is no longer needed for the project's continuation.

The bootstrap fixed point is the algorithmic instance of Lawvere's
fixed-point theorem applied to compiler construction. Lawvere's
theorem says: a self-mapping that admits an evaluation procedure
on its own representation has a fixed point. The Lean 4 compiler
maps Lean 4 source to Lean 4 binaries; the compiler's own source
is in Lean 4 (representable in its own input format); the
evaluation procedure is the bootstrap. The fixed point exists by
Lawvere; it is computed by the bootstrap procedure.

The invariant in this case is the compiler's binary. The
recompilation process is the relabeling. The invariance is that
the binary does not change after the fixed point is reached.

The compiler's binary is approximately a hundred megabytes. The
specification of the binary's behavior --- the Lean 4 language
manual, the kernel reduction rules, the elaboration algorithm, the
typeclass resolution procedure --- is hundreds of pages. The
binary implements the specification. The specification defines
the binary's behavior. The invariance is the fixed-point
relationship between the two.

This is a non-trivial property. Many implementations of a
specification do not have this property. Two C++ compilers
implementing the same C++ standard may produce different binaries
from the same source: they have different optimization strategies,
different runtime libraries, different code generators. Even if
both compilers are correct (both implement the C++ standard
faithfully), they produce different concrete outputs.

The Lean 4 compiler aims for a stronger property: the binary's
behavior is reproducible. Compiling the same source with the same
compiler produces the same binary. The binary is the canonical
implementation of the source. This canonicalness is essential for
distributed development: when two developers compile the same
source, they get the same binary, and they can be sure their
binaries behave identically.

The reproducibility is also essential for proof verification.
A Lean 4 proof is a term that the compiler checks for type
correctness. Two compilers that produce different binaries from
the same source may also check proofs slightly differently --- not
in their judgment of correct/incorrect, but in details like the
specific error messages or the timing of the check. For practical
proof verification, this is usually adequate. For theoretical
purposes (a critic might ask: did the compiler used really check
the proof?), the reproducibility provides a clean answer:
yes, this binary checked it; the binary is canonical.

The chapter's algorithmic claim is that bootstrapping is the
canonical example of admissible recompilation. The compiler is
the algorithm. Recompiling is the relabeling. The fixed-point
result is the invariance. The Lean 4 compiler is the most
visible example because the Lean 4 community has explicitly
demonstrated the bootstrap; many other compilers in the world
have implicitly bootstrapped but not always demonstrated the
fixed point publicly.

The Lean analog typeclass is \texttt{ElaborationProcess} in
\texttt{Episode12.lean}. The class formalizes the operation of
moving from load (the source) through transform (the
elaboration) to boolean (the type-check verdict). The
\texttt{HALTED} certificate marks the type-check as final. The
fixed point is achieved when the type signature of the
\texttt{ElaborationProcess} for the compiler equals the type
signature for the recompiled compiler. The fixed-point property
is checked at the type level, not at the binary level.

This is the chapter's distinguishing observation: invariance is
a type-level property, not a binary-level one. Two binaries can
differ in their byte representations but agree on their type
signatures. The agreement at the type level is what the chapter
considers invariance. The binaries are realizations; the type
signatures are the invariant content.

When the Lean compiler bootstraps, the type signatures of all
its definitions are preserved across the recompilation. The
binary may differ in compile-time metadata, in optimization
choices that produce equivalent code, in the order of unrelated
sections --- but the type signatures are identical. The
preservation of type signatures is the bootstrap's algorithmic
invariance.

The compiler community uses this invariance to certify new
versions. Before releasing a new Lean 4 compiler, the
developers verify that the bootstrap converges: the new compiler,
compiled with the old, produces a binary that, when used to
compile the new source again, produces an identical binary. The
fixed point is checked. If the check passes, the new version is
released. If it fails, the developers investigate the
discrepancy and fix it.

The fixed point is also why the compiler can be trusted to be
self-consistent. The compiler's source code includes its own
elaboration rules, its own type checker, its own kernel. Each of
these can be independently audited. Once the compiler bootstraps,
the audit certifies the entire stack: the bootstrap process
demonstrates that the auditable source produces the same binary
that is being used to build the project. The trust chain
closes.

In broader algorithmic terms, the bootstrap is the analog of a
self-referential proof. The proof says: this proof is correct.
The Godel-like self-reference is admissible because the proof
includes its own verification procedure. The compiler's source
includes its own compilation procedure; the bootstrap is the
verification.

The fixed point is the precise sense in which the compiler is
invariant under admissible relabeling. Relabeling the source's
variable names, reformatting the source's whitespace, reordering
its definitions --- any of these are admissible transformations
that leave the binary unchanged (modulo metadata). The binary
is the invariant. The source is the variable form.

The chapter's algorithmic conclusion: invariance is the fixed
point of admissible transformation. The compiler bootstraps to
its invariant binary. Any algorithm that admits relabeling has
some invariant content, and the invariant content is what
survives the relabeling. The bootstrap procedure makes the
invariance computable for the case of the compiler. For other
algorithms, the analogous procedure may be harder to specify,
but the principle holds: invariance is what does not change
under admissible relabeling.

\section{Gradient Descent Invariance}
\label{se:gradient-descent-invariance}

A neural network with 128 hidden neurons in one layer can be trained
on MNIST and reach about 98 percent accuracy. Now relabel the 128
neurons: replace neuron 1 with neuron 73, neuron 2 with neuron 41,
and so on according to some random permutation. Apply the same
permutation simultaneously to the corresponding rows of the
incoming weight matrix and columns of the outgoing weight matrix,
so the network as a whole is just the original network with its
hidden units relabeled. The loss is exactly the same as in the
original training. The accuracy is exactly the same. The weights
have been permuted consistently with the neuron relabeling, but
the network's input-output behavior is identical.

This invariance is exact only for the simultaneous permutation of
hidden units together with their adjacent weights. It does
\emph{not} say that independently retraining a relabeled network
from a fresh random initialization reaches the same loss; that is
a separate empirical claim that depends on optimization landscape
properties and is not guaranteed. The invariance the chapter
appeals to is the structural one: the loss surface is symmetric
under the action of the symmetric group $S_{128}$ on hidden
neurons.

The invariance is the chapter's question made concrete in machine
learning. The neurons are labels; the network's function is
invariant under relabeling. The labels do not affect the
network's input-output behavior; they only affect the names of
the internal computations.

The loss function is the invariant. The cross-entropy loss is a
function of the network's outputs and the true labels. The output
labels (the ten digit classes) are fixed; the loss is fixed when
the network's output distribution is fixed. The neuron labels are
internal; they do not appear in the output distribution.

This is the algorithmic content of gauge symmetry. The network has
a relabeling symmetry: any permutation of the hidden neurons,
applied simultaneously to the adjacent weights, is a symmetry of
the loss function. Gradient descent commutes with this symmetry:
running training on the relabeled-and-reweighted network is the
same as running training on the original and then applying the
relabeling. Independent retraining is a different question;
gradient descent's commutation with the simultaneous-permutation
symmetry is the claim the chapter relies on.

The Lean analog in \texttt{Episode13.lean} is the
\texttt{Measurement.le} relation. The relation compares two
measurements (origin, distance, speed) as ordered stages. The
order is the relabeling-invariant content: regardless of how the
stages are named, the relation between them is the same.
\texttt{LeanProcess} carries the algorithmic content;
\texttt{MEASURED} certifies that the result is invariant under
the typeclass-instance relabelings.

The invariance has consequences for training. Different
initialization seeds produce different specific weights, but the
trained networks (after convergence) have the same loss. The loss
landscape's minima are related by the relabeling group; the
training process finds one of them, and the choice of which one
depends on the seed. The function the network implements is the
invariant; the specific weights are the variable form.

This is also why deep learning works in practice. The network has
enormous flexibility in its specific weights, but the function it
implements is constrained by the loss function. The invariance is
the structure that makes the optimization tractable: gradient
descent on a loss function with a vast symmetry group finds any
of the symmetry-equivalent minima. The found minimum represents
the same function, regardless of which specific representation
was reached.

\section{Ising Model Phase Transition}
\label{se:ising-model-phase-transition}

The two-dimensional Ising model has a square lattice of spins, each
of which can take one of two values: up or down. Adjacent spins
interact: aligned spins (both up or both down) contribute a small
negative energy; anti-aligned spins contribute a small positive
energy. At temperature $T$, the system explores configurations
according to the Boltzmann distribution: lower-energy configurations
are more probable.

At high temperatures, the thermal energy overwhelms the spin
interactions; the spins fluctuate randomly; the average
magnetization (the fraction of up spins minus the fraction of down
spins) is zero. At low temperatures, the thermal energy is small
compared to the spin interactions; the spins prefer to align; the
system has a nonzero magnetization, either positive (mostly up) or
negative (mostly down). Between high and low temperatures is a
critical temperature $T_c \approx 2.269$ (in units where the
interaction is 1) at which the system undergoes a phase transition.

The Hamiltonian has a $\mathbb{Z}_2$ symmetry: flipping every spin
(up to down, down to up) leaves the energy unchanged. The
configuration ``mostly up'' has the same energy as the configuration
``mostly down.'' Above $T_c$, the system has the symmetry: the
average magnetization is zero, consistent with the symmetry. Below
$T_c$, the system spontaneously breaks the symmetry: the magnetization
is either positive or negative, but not zero. Which sign is chosen
depends on initial conditions and small fluctuations; the choice is
random.

The order parameter is the invariant. The order parameter is the
magnetization. Above $T_c$, it is zero; the symmetry is unbroken.
Below $T_c$, it is nonzero; the symmetry is broken. The order
parameter records the breaking. It is the invariant of the system's
state under the symmetry's possible orbit: above $T_c$, the orbit
is a single point (zero); below, the orbit has two points (positive
and negative).

The Metropolis algorithm simulates the Ising model. The algorithm
flips spins randomly, accepting flips that lower energy and
sometimes accepting flips that raise energy (with probability
$e^{-\Delta E/T}$). Over time, the algorithm explores configurations
according to the Boltzmann distribution. After many steps, the
algorithm's average magnetization matches the equilibrium
magnetization.

The Lean analog is the typeclass cascade for phase transitions in
\texttt{Measurement.le}. A measurement that crosses a phase
transition (a structural discontinuity in the system's behavior)
is recorded as a discrete event: the order parameter is one value
before the transition and another after. The transition itself is
an admissible refinement step that the algorithm tracks
explicitly.

The chapter's claim: the invariant in a symmetry-broken phase is
the order parameter. The symmetry has multiple admissible orbits
(in the broken phase); each orbit's representative is the order
parameter's specific value. The algorithm carries the order
parameter as the invariant that distinguishes the broken-symmetry
orbits from each other.

\section{Feynman Diagram as Computation}
\label{se:feynman-diagram-as-computation}

In quantum electrodynamics, the electron's anomalous magnetic
moment (the deviation of its magnetic dipole from the Dirac value)
is one of the most precisely measured quantities in physics. The
theoretical prediction is computed using Feynman diagrams: a
graphical representation of the perturbative expansion of the
quantum amplitude.

The simplest diagram contributing to the anomalous moment is the
one-loop vertex correction: an electron scatters off an external
photon, but during the scattering, the electron emits and reabsorbs
a virtual photon. The diagram is a triangle: two electron lines
joining at the external photon vertex, with the virtual photon
running between them.

The diagram's mathematical content is an integral over the
internal momentum of the virtual photon:
\[
\Gamma_\mu \propto \int \frac{d^4 k}{(2\pi)^4} \, \frac{1}{k^2 + i\epsilon}
\cdot (\text{spinor structure}).
\]
The integral diverges (it is ultraviolet divergent), and the
regularization and renormalization procedures handle the
divergence. The renormalized amplitude is finite and gives the
specific contribution to the anomalous moment.

The diagram is invariant under relabeling of the internal momentum.
If we substitute $k \to k + p$ for any constant momentum $p$, the
integral's value does not change (after appropriate
renormalization). The internal momentum is a dummy variable. The
diagram's value depends on the external momenta (the photon's
momentum, the electron's momenta in and out), not on how we choose
to label the internal momentum.

This relabeling invariance is the chapter's question in particle
physics. The Feynman diagram is the algorithm. The internal
momentum is the variable; the external momenta are the inputs. The
amplitude is the output. Relabeling the variable does not change
the output. The output is the invariant content.

The Feynman rules formalize this. The rules specify, for each
particle and each interaction, how to translate a diagram into an
integral. The rules are invariant under the relabeling: applying
the rules to a diagram and to its relabeled version produces the
same integral, evaluated with different dummy variable names. The
integral is the algorithmic content; the labels are operational.

In Lean, the analog is the alpha-equivalence of bound variables.
Two lambda terms that differ only in the names of bound variables
are alpha-equivalent: \texttt{$\lambda$ x => x + 1} and
\texttt{$\lambda$ y => y + 1} are the same function. The compiler
treats them as identical for purposes of type-checking and
evaluation. The variable name is the label; the function is the
invariant.

The chapter's algorithmic claim: alpha-equivalence is the universal
form of relabeling invariance. Every algorithm that uses bound
variables has this invariance: the algorithm's output depends on
the variables' types and roles, not on their specific names. The
chapter's typeclasses like \texttt{ElaborationProcess} respect this
invariance: the elaborator does not care whether a definition
names its argument \texttt{x} or \texttt{n}.

\section{ElaborationProcess in Lean}
\label{se:elaborationprocess-in-lean}

\texttt{Episode12.lean} defines \texttt{ElaborationProcess} as the
typeclass that captures the stages of compilation. The structure
has the shape:

\begin{leanexcerpt}{ElaborationProcess}
structure ElaborationProcess
    (Value : Type)
    (Carrier : CarrierProcess Value)
    [DISTINGUISHABLE Value Carrier] ... [LOGICAL Value Carrier]
  where
  stamina        : HeartbeatProcess Value Carrier
  calibration    : Calibration.EKG
  computer_state : ComputerProgram
  execute        : ComputerProgram $\to$ ComputerProgram := ...
\end{leanexcerpt}

The structure carries the full project cascade through
\texttt{LOGICAL}, plus four fields: \texttt{stamina} (the
underlying \texttt{HeartbeatProcess}), \texttt{calibration} (the
\texttt{Calibration.EKG} reference), \texttt{computer\_state} (the
current \texttt{ComputerProgram}), and the function
\texttt{execute} that advances a program through one stage. The
function's three cases correspond to the three constructors of
\texttt{ComputerProgram}: \texttt{load} (initialize the
elaboration state from the source), \texttt{transform} (apply
unification, typeclass resolution, and other elaboration steps),
and \texttt{boolean} (produce the final type-check verdict).

The accompanying typeclass \texttt{HALTED} certifies that the
\texttt{boolean} stage has been reached:

\begin{leanexcerpt}{HALTED}
class HALTED
    (Value : Type)
    (Carrier : CarrierProcess Value)
    [DISTINGUISHABLE Value Carrier] ... [LOGICAL Value Carrier]
  where
  scientific_paper : ElaborationProcess Value Carrier
  halted? : ComputerProgram $\to$ ComputerProgram $\to$ Prop := ...
\end{leanexcerpt}

A \texttt{HALTED} instance carries the cascade through
\texttt{LOGICAL}, a reference \texttt{scientific\_paper} (the
underlying \texttt{ElaborationProcess}), and a binary decision rule
\texttt{halted?} on \texttt{ComputerProgram}s. The decision rule
returns \texttt{True} exactly when the second program is a
\texttt{.boolean} constructor, certifying that the elaboration has
reached its boolean stage. Only \texttt{HALTED} elaborations are
admissible as final type-checked declarations. The \texttt{load}
and \texttt{transform} cases of \texttt{ComputerProgram} are
intermediate; they return \texttt{False} from \texttt{halted?}.

The invariant under recompilation is the result of the
\texttt{boolean} stage. Two compilations of the same source
produce the same \texttt{HALTED} result. The \texttt{load} and
\texttt{transform} stages may differ in their internal details
(the specific order of unification steps, the specific typeclass
instances selected, the specific cache states), but the
\texttt{boolean} verdict is the same.

This is the chapter's invariance principle applied to compilation.
The compilation process has many internal details that vary across
runs. The final type-check result does not. The result is the
invariant; the internal details are the variable form.

The \texttt{CompilerTape}, \texttt{CompilerOutput}, and
\texttt{COMPILED} typeclasses in \texttt{Episode14.lean} formalize
the staged artifact of compilation. The tape is the compilation
trace; the output is the final binary; \texttt{COMPILED} certifies
that the binary is the canonical output for the source. The
canonical-ness is the bootstrap fixed point: the output is
invariant under admissible recompilation.

The \texttt{Measurement.le} relation in \texttt{Episode13.lean}
extends the invariance to staged measurements. A measurement has
three stages: origin (where the measurement begins), distance
(how far the measurement extends), and speed (how fast the
measurement evolves). The order on measurements is the
relabeling-invariant content: one measurement refines another
when its origin, distance, and speed all refine.

The \texttt{LeanProcess} typeclass marks a process as
implementing the Lean elaboration discipline. A process is
\texttt{LeanProcess} when it preserves the invariants of
elaboration across its internal operations. The
\texttt{MEASURED} certificate is the formal proof that the
process's output is invariant under the admissible
relabelings.

The chapter's closing observation: the invariance of computation
under admissible relabeling is what makes the compiler
trustworthy. If the compiler's output depended on internal
details that vary across runs, the binary would not be a
reliable artifact. The compiler's invariance is what enables
reproducible builds, deterministic deployments, and verifiable
proof checks. The chapter's typeclass machinery is the formal
discipline that enforces the invariance.

\begin{bridge}

The chapter's question was what remains invariant under admissible
relabeling or recompilation.

The answer is the type-level content. The compiler bootstrap produces
a fixed-point binary whose type signatures are invariant under
recompilation. Gradient descent training produces a network whose
loss is invariant under hidden neuron relabeling. The Ising model
below the critical temperature has spontaneously broken symmetry
whose order parameter is the invariant. Feynman diagram amplitudes
are invariant under relabeling of internal momentum. Lean's
\texttt{ElaborationProcess} preserves the type-check verdict
across the internal details of elaboration; \texttt{Measurement.le}
preserves the staged content of measurements.

What remains is the final chapter: append-only computation,
incompressibility, the type-check that closes the book. The
algorithm's record grows monotonically; the growth cannot be undone;
the final type-check is the elaborator's verdict on whether the
accumulated proof burden has reached a consistent state.

\end{bridge}

\begin{coda}{Theater Production with Rotating Understudies}

The theater is in its fifteenth year of running the same play. The
play is a two-act drama about a family across three generations,
written by a deceased playwright, directed initially by a young
director whose career has since taken him elsewhere. The current
director has been with the production for nine years; he replaced
the original director, who replaced the founding director, who
left to direct a film.

The cast has rotated entirely. None of the original eight actors
remain. The current cast is the third generation of actors to play
these roles. The actor playing the grandfather has been in the
production for four years; he replaced the previous grandfather,
who left to do a film of his own. The actor playing the mother is
new this season. The understudies for the three principal roles
have themselves rotated: the current understudies are the second
generation of understudies in their respective roles.

The set has been rebuilt twice. The original set was designed by
the founding director's frequent collaborator; that set was retired
after eight years, when the wood structure had developed cracks
that no longer matched the original visual intention. The
replacement set was designed by a different scenographer; it is
the second set, currently in use. The props have been replaced
piecemeal as they have worn out: the kitchen table is the third
table since the production opened; the family photograph in the
center is the fifth print (digital files restored from the
original photographer's archive); the rocking chair is the fourth.

The lighting rig has been replaced entirely. The original rig used
incandescent bulbs and required two stagehands to operate. The
current rig is LED, programmed by a single technician at a digital
console. The lighting cues are different now (the LED rig allows
faster transitions than the original incandescents), but the
lighting design's overall effect is the same: each scene has the
intended mood, the actors are illuminated correctly, the audience's
attention is directed appropriately.

And yet the production is ``the same play.'' Three generations of
actors, two sets, three props sets, two lighting rigs, three
directors --- and the play itself, as performed nightly, is
recognizably the same work. Audiences who saw the production fifteen
years ago and return today recognize the play. The critics
describe it as the same play. The publishing house lists it as
the same play. What survives?

Not the actors, the sets, the lights, the directors. Not any
particular performance. What survives is the relational structure
of the performance: the script (the words spoken), the blocking
(the spatial pattern of movement on stage), the timing (the rhythm
of the scenes), the emotional arc (the trajectory of the family's
story across the play's three generations).

The script is the most obviously invariant. The words are the same
words. The same lines have been spoken, by three different casts,
in the same order, with the same pauses, for fifteen years. The
script is the binary of the production: the canonical output that
all the variables (cast, set, lights) realize.

The blocking is invariant up to small variations. Each director has
adjusted the blocking slightly --- different placements of the
characters, different patterns of movement --- but the overall
structure is preserved. A reviewer who saw the production in year
one and saw it in year fifteen would describe the blocking as
``the same.'' The variations are within an admissible equivalence
class.

The timing is invariant up to small variations. Different casts
have different rhythms, but the play's overall pacing is preserved:
the first act runs eighty minutes, the intermission fifteen, the
second act ninety. The total running time has varied by only a
few minutes over the years.

The emotional arc is invariant. The audience experiences the same
journey across the play's three generations. Joy in the first act,
tension building in the middle, catharsis in the third. The arc is
the play's deepest invariant: it survives the changes in cast,
set, lighting, and direction.

What changes is the implementation. The implementation is the
specific actors, specific costumes, specific scenery, specific
lighting cues. The implementation is the variable form. The
implementation can be replaced; the play remains.

This is the chapter's structure in theatrical form. The play is
the type --- the abstract specification of the performance. The
production is the binary --- the realized instance. Different
binaries (different casts and crews) implement the same type. The
implementations differ in details; the type is invariant.

The director's job, when bringing on a new cast member, is to
maintain the invariance. The new actor learns the lines (the
script invariant), learns the blocking (the spatial invariant),
learns the timing (the temporal invariant), and integrates with
the existing cast (the relational invariant). The director
shapes the new performance to be within the equivalence class of
admissible realizations of the play.

When the director makes a change --- a new lighting cue, a slight
adjustment to the blocking, a different prop --- the change is
admissible if it remains within the equivalence class. If the
change leaves the class (a major rewrite of a scene, a different
ending, a different setting), the production becomes a different
play: a new type, not a new instance of the same type. The
boundary of the equivalence class is the play's identity.

The chapter's algorithmic claim: this is the universal pattern of
invariance under admissible relabeling. The play is the
\texttt{COMPILED} artifact. The actors are the
\texttt{Measurement} stages: each cast is one realization. The
director is the \texttt{ElaborationProcess} that produces the
realization from the script. The \texttt{HALTED} certificate is
the curtain call: the performance has reached its conclusion. The
play, as carried across all the changes, is the
\texttt{COMPILED} invariant.

The producer is reading the audience reviews from this week. The
reviews are positive. The play is in its fifteenth year, and
audiences are still finding it powerful. The cast has rotated, the
set has been rebuilt, the lights are LED, the director is the
third in the line --- and the play is the same play.

The next season's plans include an out-of-town tour: the production
will travel to four cities. New venues, possibly modified blocking
to fit the different stages, the same script, the same actors. The
tour is another set of admissible recompilations. The play remains.

\end{coda}
